\chapter{Données et limites}
\label{chap:donnees}

% MIGRÉ AUTOMATIQUEMENT depuis memoire/main.tex (plages 1002-1433).
% Hiérarchie descendue d'un cran, labels préfixés "don:".
% RELECTURE FAITE (août 2026). Aucune phrase du chapitre n'annonce la copule ni
% la variable latente comme LE modèle. Ajouts d'août : la table des trois
% fréquences, le seuil de collecte, et la section des deux biais chiffrés
% (élasticité de taille corrigée par troncature, chain-ladder de remontée).
%==============================================================

La modélisation ne repose sur aucune base unique. Chaque source mobilisée a son périmètre,
son biais et sa granularité, et c'est leur articulation, non le choix d'une source
privilégiée, qui fait la démarche.

%--------------------------------------------------------------
\section{Cartographie des sources et de leurs rôles}
\label{sec:carto}
%--------------------------------------------------------------

Aucune base n'est utilisée seule : chaque source alimente une composante précise du modèle,
selon ce qu'elle sait porter.

\begin{center}
\renewcommand{\arraystretch}{1.3}
\begin{tabular}{>{\raggedright}p{3cm} >{\raggedright}p{5.5cm} >{\raggedright\arraybackslash}p{5.5cm}}
\toprule
\rowcolor{lightblue}
\textbf{Source} & \textbf{Rôle dans le modèle} & \textbf{Justification / limite} \\
\midrule
PRC 2025 & Sévérité alternative via conversion Jacobs ; fréquence de référence $\lambda_{ref}$, sans rôle de niveau &
Périmètre défini, mais sévérité dérivée d'un modèle log-log \\
\addlinespace
SAS OpRisk Global & Sévérité (montants réels) et fréquence d'entrée du moteur ($\lambda$) &
Montants observés, mais biais vers les grandes entités financières \\
\addlinespace
Hackmageddon & Parts par vecteur d'attaque (structure du risque) &
Ne fixe aucun niveau : pondère le multiplicateur qui fait passer $\lambda$ d'un état de conformité à l'autre \\
\bottomrule
\end{tabular}
\end{center}

Cette répartition écarte un écueil fréquent : traiter une base généraliste comme si elle
pouvait, seule, calibrer un modèle fréquence-sévérité. Chaque source est mobilisée là où elle
est la plus fiable, et ses limites sont déclarées.

%--------------------------------------------------------------

%--------------------------------------------------------------

L'usage d'un modèle interne, même partiel, est encadré par la directive Solvabilité~II.
L'Article~19 du Règlement Délégué (UE) 2015/35 impose que les données servant au calcul de la
distribution de probabilité prévisionnelle respectent trois critères cumulatifs : exactitude,
exhaustivité et pertinence.

Sur le risque cyber, aucune base publique ou commerciale ne satisfait ces trois critères.
L'actuaire doit donc documenter les biais des données historiques et justifier les
retraitements. La validité du capital calculé ne tient pas à la perfection de la donnée, qui
n'existe pas, mais à la déclaration des limites et à la robustesse du modèle face à ces
biais.

%--------------------------------------------------------------

%--------------------------------------------------------------

% SOURCES-SCRIPTS: 63

\section{Trois biais statistiques à déclarer avant tout usage}

L'utilisation conjointe des bases \textit{Privacy Rights Clearinghouse} (PRC) et SAS OpRisk Global permet de croiser les perspectives, mais nécessite d'identifier trois biais statistiques majeurs qui altèrent l'estimation de l'indice de queue $\xi$.

\paragraph{1. Le biais de sélection (Selection Bias).}
Ce biais affecte particulièrement la base OpRisk Global, qui se nourrit principalement d'événements rendus publics dans la presse financière ou les rapports annuels. Par construction, un incident cyber n'atteint ce seuil de visibilité que s'il frappe une entité d'une certaine taille (généralement des institutions systémiques ou de grandes banques internationales) ou si son montant est exceptionnellement élevé.
Ce biais de taille sur-représente les événements extrêmes et sous-représente la sinistralité
d'attrition. Calibrer directement sur cette base un modèle destiné à une entité de taille
moyenne, une mutuelle régionale par exemple, surestime mécaniquement la probabilité de perte
extrême. C'est pourquoi la fréquence et la sévérité sont ramenées à la taille de l'entité par
deux élasticités estimées (\S\ref{sec:biais-chiffres} et chapitre~\ref{chap:resultats}).

\paragraph{2. Le biais de déclaration (Reporting Bias).}
La base PRC est soumise à un biais de déclaration lié aux asymétries réglementaires. Aux États-Unis, la notification des fuites de données personnelles (Data Breaches) est obligatoire dans la plupart des États, ce qui rend la base PRC très riche sur ce vecteur d'attaque. En revanche, un rançongiciel (Ransomware) qui paralyse les serveurs de l'entreprise mais n'exfiltre aucune donnée personnelle peut ne déclencher aucune obligation de notification légale (avant l'entrée en vigueur d'actes récents comme le CIRCIA).
Calibrer une fréquence globale sur la PRC sous-estime donc le risque d'interruption
d'activité au profit du seul risque de compromission de données. C'est pourquoi la fréquence
qui porte le capital n'est pas lue sur la PRC mais sur OpRisk, qui recense les pertes matérielles
quel que soit le vecteur, et pourquoi la répartition par vecteur, rançongiciels compris, est prise
dans la base Hackmageddon plutôt que dans la PRC.

\paragraph{3. Le biais de survie (Survivorship Bias).}
Dans les bases de pertes opérationnelles, la sévérité maximale observée est paradoxalement bornée par la capacité de survie de l'entité. Une cyberattaque ayant entraîné la faillite pure et simple de l'entreprise (ruine) disparaît souvent des bases de pertes opérationnelles classiques, car il n'y a plus d'entité survivante pour reporter la perte financière finale et les coûts de remédiation à long terme. Ce biais tronque artificiellement l'extrême queue de distribution. La théorie des valeurs
extrêmes ne le corrige pas : une loi de Pareto généralisée ajustée sur un échantillon tronqué
hérite de la troncature, et si elle permet d'extrapoler au-delà du maximum observé, elle le fait
à partir d'une queue allégée. Le biais de survie va donc dans le sens d'une sous-estimation de la
queue, et il s'ajoute aux écarts anti-conservateurs déclarés au chapitre~\ref{ch:robustesse}.

%--------------------------------------------------------------

%--------------------------------------------------------------

\section{L'hétérogénéité temporelle et l'actualisation des sinistres}

L'hétérogénéité temporelle des bases de données constitue une autre infraction potentielle au critère d'exactitude de Solvabilité II. La base OpRisk Global agrège des pertes de $2000$ à $2026$. Employer un montant de $2018$ tel
quel dans un modèle de capital de $2026$ fausse le résultat, l'inflation des sinistres
technologiques ne suivant pas l'indice des prix à la consommation.

Le coût d'un sinistre cyber est principalement tiré par trois facteurs :
\begin{itemize}
    \item \textbf{L'inflation salariale IT :} Le coût des experts en remédiation (forensic, restauration de systèmes) connaît une inflation nettement supérieure à l'indice des prix à la consommation, en raison de la pénurie de talents en cybersécurité.
    \item \textbf{L'inflation juridique et réglementaire :} L'entrée en vigueur du RGPD en 2018 a drastiquement modifié l'échelle des sanctions. Une fuite de données en 2015 coûtait infiniment moins cher en amendes administratives que la même fuite en 2025.
    \item \textbf{L'escalade des rançons :} La professionnalisation du cybercrime (Ransomware-as-a-Service) a transformé les demandes de rançons de quelques milliers de dollars en 2016 à plusieurs dizaines de millions d'euros en 2026.
\end{itemize}

Pour respecter les exigences de l'Article 19, les sinistres historiques doivent être actualisés en euros de l'année de modélisation (2026). Soit $L_t$ la perte observée l'année $t$, la perte actualisée $\tilde{L}_{2026}$ s'écrit :
\begin{equation}
\tilde{L}_{2026} = L_t \times \prod_{k=t}^{2025} (1 + i_k^{cyber})
\end{equation}
où $i_k^{cyber}$ est un taux d'inflation propre aux sinistres technologiques, que les indices
sectoriels situent entre $5$ et $10\,\%$ par an. La base OpRisk porte deux montants pour chaque
perte : le montant enregistré l'année de l'événement et sa valeur courante, corrigée de
l'inflation générale. La chaîne publiée lit le premier, donc des montants \emph{nominaux}. Le
rapport entre les deux colonnes vaut un pour les années récentes et $1{,}747$ en $2004$, ce qui
fixe l'inflation implicite à $2{,}19\,\%$ par an. Cette homogénéisation n'est pas appliquée à la
calibration gelée : elle minore les sinistres anciens ramenés en euros courants, donc la queue et
le capital, et son effet sur la dérive des pertes est mesuré au \S\ref{soc:sec:backtest}. Aucun
indice d'inflation spécifique au cyber ne fait aujourd'hui consensus, de sorte que la valeur
courante de la base, fondée sur l'indice général, reste elle-même un minorant de l'actualisation
que l'équation demande.
% SOURCES-SCRIPTS: 102

Par ailleurs, le traitement des données manquantes dans la PRC (qui recense parfois le nombre d'enregistrements compromis sans le montant financier) est géré par la conversion de \citet{jacobs2014}. Bien qu'utile comme \textit{proxy}, ce modèle de régression log-log lisse la variance et minore les extrêmes. C'est le motif de
ne pas plafonner la calibration issue directement d'OpRisk, où les montants sont observés de
bout en bout.

%--------------------------------------------------------------

%--------------------------------------------------------------

Les limites épistémiques et statistiques exposées ci-dessus contraignent aujourd'hui le calcul du SCR cyber à de larges intervalles d'incertitude. Cependant, cette situation est transitoire. L'apport fondamental du règlement DORA ne réside pas uniquement dans l'imposition de normes de sécurité (Chapitre II), mais dans la création d'un cadre de notification harmonisée des incidents majeurs liés aux TIC (Chapitre III, Article 19 de DORA).

DORA impose aux entités financières de reporter à leur autorité compétente (l'ACPR en France) les incidents majeurs selon une taxonomie standardisée (causes, vecteurs, temps de résolution, pertes financières directes et indirectes). À terme, le règlement prévoit la centralisation de ces rapports au sein d'un pôle de l'Union européenne (EU Hub).

Ce registre atténuerait les trois failles actuelles sans les faire disparaître :
\begin{itemize}
    \item \textbf{Biais de déclaration :} la notification étant obligatoire, la base couvrirait
      les incidents \emph{majeurs} au sens du règlement. Elle resterait donc tronquée au seuil
      de matérialité, et non exhaustive : le biais change de nature plutôt qu'il ne disparaît.
    \item \textbf{Biais de sélection :} le registre couvrirait toutes les strates du marché
      financier, courtiers, mutuelles et sociétés de gestion compris, ce qui donnerait la
      granularité nécessaire pour conditionner la sévérité par taille d'entité.
    \item \textbf{Visibilité sur les tiers :} le registre cartographierait l'implication des
      prestataires tiers TIC. C'est exactement la donnée qui manque pour identifier la direction
      de la contagion, et le chapitre~\ref{chap:partielle} chiffre en euros de capital ce que
      son obtention ferait gagner.
\end{itemize}

L'architecture retenue ici (approche par distribution des pertes, théorie des valeurs
extrêmes, facteur de Vasicek, allocation d'Euler) ne dépend pas de la source : le modèle et les
scripts sont agnostiques aux données d'entrée. Substituer une base d'incidents DORA harmonisée
aux bases PRC et OpRisk ne demanderait donc aucune réécriture. Le gain porterait sur la qualité
des entrées, non sur la structure du modèle, et il resterait borné par ce que le registre ne
couvrira pas : les incidents sous le seuil de matérialité, et la direction de la contagion tant
que le registre ne l'horodate pas.

\paragraph{Automatisation du reclassement par NLP.} En attendant ce registre harmonisé, une
partie du travail de fiabilisation peut déjà être outillée. Le reclassement manuel par grille
de mots-clés opéré plus haut sur la taxonomie Hackmageddon, nécessaire pour isoler le
ransomware qui n'existait pas comme catégorie en 2023, illustre une limite récurrente des
bases de sinistralité cyber : une taxonomie déclarative instable, dont la correction repose
sur un travail manuel documenté mais non reproductible à grande échelle.

Les travaux de recherche menés chez Nexialog Consulting sur le traitement automatique du
langage appliqué aux bases de sinistralité cyber (rapports d'incidents, bases de
vulnérabilités, publications techniques) ouvrent une voie pour systématiser ce reclassement.
Le gain serait double : une classification supervisée validée statistiquement remplacerait la
grille ad hoc, et la procédure deviendrait transposable à d'autres bases, OpRisk comme le
futur registre DORA. Faute de corpus annoté dans le périmètre de ce travail, cette piste n'a
pas été mise en \oe uvre ; elle constitue un prolongement direct.

%--------------------------------------------------------------

%--------------------------------------------------------------

La base \emph{Privacy Rights Clearinghouse} (PRC) fournit la fréquence de
référence $\lambda_{ref}$, qui compte des notifications sans seuil de
matérialité et ne porte donc aucun niveau de capital (\S suivante). Son intérêt
tient à son périmètre défini et à sa couverture longitudinale.

La PRC joue toutefois un second rôle, plus fragile : celui de source de
sévérité indirecte, via une conversion dite de Jacobs, qui traduit un nombre
d'enregistrements compromis en un montant de perte financière à partir d'un
modèle log-log. Ce mécanisme permet d'élargir le spectre de tailles
d'incidents disponibles, mais introduit un biais de modèle supplémentaire :
la sévérité qui en résulte n'est pas observée directement, elle est
\emph{dérivée}. La calibration en queue de distribution obtenue sur la PRC,
réalisée en propre sur les 15\,053 incidents 2019--2025 de la base brute,
via la conversion Jacobs $\ln(L) = 7{,}68 + 0{,}76\,\ln(X)$, donne un
paramètre de forme $\hat{\xi}_{\text{PRC}} = 1{,}03$ (IC90\%
$[0{,}97~;~1{,}09]$), caractéristique d'une queue extrêmement lourde, à
espérance infinie, résultat qui, pris seul, rendrait toute simulation
numériquement instable et qui doit donc être lu avec le plafond de sévérité
présenté au \S\ref{sec:portee-donnees}.

La PRC reste utile malgré ces limites parce qu'elle est la seule source
alternative de sévérité : elle éprouve la calibration OpRisk sur un autre
périmètre, sans porter elle-même le niveau.

% SOURCES-SCRIPTS: 57 62 63

\section{Les trois fréquences, et le seuil de collecte qui les sépare}

% SOURCES-SCRIPTS: 08b 60 62

\paragraph{Le seuil de collecte et la double lecture de $\lambda$.}
La fréquence est la seule brique du modèle dont la valeur dépende
explicitement du \emph{périmètre de comptage} de la source, et cette dépendance
doit être déclarée avant tout usage du niveau. Un contrôle mené sur le panel
firme-année du SAS OpRisk (sous-catégories TIC, 2005--2022, fenêtre
d'observation propre à chaque firme, de sorte qu'une année sans événement soit
un vrai zéro) donne $0{,}10$ événement par firme et par an toutes firmes
confondues, et $0{,}21$ pour les firmes à au moins dix événements. Ces deux
valeurs sont des \emph{seaux} et ne décrivent aucune entité en particulier : le
chapitre~\ref{chap:resultats} estime $\lambda$ comme une fonction de la taille de
firme et la lit à celle de l'entité visée, ce qui donne $0{,}092$. Toutes trois
sont sans commune mesure avec le taux posé dans le modèle. L'écart n'est pas une erreur de
calage : la base OpRisk ne retient que les pertes au-dessus d'un seuil
de collecte, si bien que le taux qu'on y lit est celui des événements TIC
\emph{matériels}, générateurs d'une perte enregistrée, et non celui de tous les
incidents au sens du reporting DORA. C'est la maille pertinente pour un modèle
de capital, mais elle n'est pas la maille réglementaire, et les deux ne doivent
jamais être comparées terme à terme.

Il en découle une règle de lecture qui vaut pour tout le mémoire. Le taux
$\lambda$ est une propriété d'entité, lisible sur un panel
firme-année ; la charge systémique commune, elle, ne se lit qu'au niveau
\emph{agrégé} et après retrait de tendance, car la surdispersion d'un panel
firme-année est dominée par l'hétérogénéité de taille des firmes, non par un
choc commun. Confondre les deux niveaux d'observation reviendrait à imputer au
risque systémique ce qui n'est qu'un effet de taille. Cette séparation
commande l'interprétation des niveaux de capital du
chapitre~\ref{chap:resultats} : selon l'échelle à laquelle $\lambda$ est lu,
le même modèle décrit un secteur ou une entité, et la
table~\ref{tab:echelle} en donne la traduction chiffrée.

\begin{attention}
\textbf{Trois nombres portent le nom de \og{}fréquence\fg{}, et il faut les distinguer une
fois pour toutes.} Le mémoire mobilise trois taux, à trois niveaux d'observation. Les
confondre serait une faute de lecture, et l'on prend donc soin de les nommer séparément.
\begin{center}\small
\renewcommand{\arraystretch}{1.3}
\begin{tabular}{@{}l l r p{5.1cm}@{}}
\toprule
\textbf{Symbole} & \textbf{Source} & \textbf{Valeur} & \textbf{Ce qu'il compte} \\
\midrule
$\lambda_{ref}$ & PRC & $341$/an & notifications de violation de données des
\emph{organisations financières} de la base PRC (code \texttt{BSF}), et non de la base
entière, qui en compte $2\,150$ par an. Sert à \emph{répartir} par vecteur d'attaque, pas à
fixer un niveau de capital. \\
$\lambda$ & OpRisk & $21{,}6$/an & incidents \textsc{tic} \emph{matériels} du secteur
financier mondial. C'est l'entrée du moteur de cascade et le seul des trois qui
porte le capital. \\
$\lambda_{\text{entité}}$ & OpRisk & $0{,}092$/an & les mêmes incidents matériels, rapportés
à \emph{une} firme et lus \emph{à la taille} de l'entité visée, non dans un seau. C'est
l'échelle de lecture d'une entité soumise à DORA. \\
\bottomrule
\end{tabular}
\end{center}
\vspace{2pt}
Le facteur $3\,719$ entre les extrêmes n'est donc pas une incohérence : c'est l'écart entre
un comptage de notifications, un comptage de pertes matérielles sectorielles et un comptage
de pertes matérielles par entité. Il est recalculé à chaque exécution du pipeline, de sorte
qu'il ne peut plus se désynchroniser des trois fréquences dont il découle.
% SOURCES-SCRIPTS: 67
\textbf{Seul $\lambda$ alimente le
moteur}, et c'est lui
seul qui doit être invoqué quand le mémoire parle de \og{}la fréquence du modèle\fg{}.
\end{attention}

\section{Deux biais déclarés et leur effet chiffré sur le résultat}
\label{sec:biais-chiffres}

Déclarer un biais ne coûte rien ; le mesurer engage. Les deux limites les plus souvent
opposées à une calibration sur OpRisk, le biais de taille et la troncature à droite, sont
ici chiffrées (figure~\ref{fig:biais}).
% SOURCES-SCRIPTS: 57

\paragraph{Les deux mesures.} Le biais de taille est réel mais faible : corrigée de la troncature
au seuil de collecte, l'élasticité de la sévérité à la taille vaut $b=0{,}087$, et la sévérité
transposée à l'entité notionnelle n'est multipliée que par $0{,}85$. La troncature à droite se
corrige par chain-ladder : les $97$ événements observés sur 2023--2025 deviennent $255$ à
l'ultime, et le risque \textsc{tic} n'a pas baissé après 2022, il n'est pas encore remonté. Les
régressions, le triangle de développement et la figure sont au \S\ref{sec:ann-biais-chiffres}.
% SOURCES-SCRIPTS: 67

\begin{cle}
\textbf{Le résultat qui compte, et qui n'était pas prévisible.} Les deux corrections vont en
sens opposés et se compensent presque exactement sur la perte moyenne :
$0{,}861\times1{,}147=0{,}987$, soit un effet net de $1{,}3\,\%$, inférieur à l'incertitude
d'estimation de chacune prise isolément. Elles ne se compensent pas sur le capital.
À queue lourde, la VaR $99{,}5\,\%$ obéit au principe du grand saut unique : elle répond à la
\emph{sévérité} bien plus qu'à la \emph{fréquence}. La correction de taille agit sur la
sévérité et se transmet donc au quantile ; celle de troncature agit sur la fréquence et s'y
transmet mal. Deux biais qui s'annulent sur l'espérance peuvent ainsi ne pas s'annuler sur le
quantile, et \emph{c'est le quantile qui fait le SCR}. Une compensation constatée en moyenne
n'autorise donc pas à ignorer les deux biais.
\end{cle}

%--------------------------------------------------------------

%--------------------------------------------------------------

\section{Hackmageddon : structure du risque, jamais niveau}
\label{sec:hackmageddon}

Cette source a déjà été présentée au \S\ref{sec:carto} comme celle qui porte la
\emph{structure} du risque et non son niveau. Reste à en donner le périmètre exact et les
deux limites qui bornent son usage.

\paragraph{Périmètre consolidé.} Le décompte retenu porte sur le semestre complet et contigu
janvier--juin 2026, soit 1\,041 incidents. La couverture pluriannuelle initialement
envisagée a été abandonnée : l'année 2024 est absente du jeu consulté, et 2025 n'est
disponible que partiellement. La base est donc exploitée pour la \emph{structure} du risque,
non pour son \emph{niveau} ou sa dynamique temporelle : elle ne fixe aucune fréquence, et ses
parts ne servent qu'à pondérer le multiplicateur de fréquence entre états
(\S\ref{sec:quatre-canaux}).

\begin{attention}
\textbf{Ces deux décomptes sont une citation, non un résultat du dispositif.} Contrairement
aux bases PRC et OpRisk, versionnées dans le dépôt du mémoire et rejouables par script, le
jeu Hackmageddon a été consulté puis non conservé. Les $1\,041$ incidents et les
$840$ à vecteur identifié ne sont donc reproductibles par aucun script, et le contrôle de
non-régression du dépôt les signale explicitement comme tels plutôt que de les compter parmi
les sorties vérifiées. Ils sont à lire comme un chiffre de rapport de place, avec sa date de
consultation, et non comme une sortie du modèle. Ces décomptes fixent les parts $\pi_v$ de
répartition par vecteur, qui ne portent aucun niveau mais pondèrent le multiplicateur de
fréquence entre états de conformité (\S\ref{sec:quatre-canaux}) : leur statut de citation
s'étend donc à ce multiplicateur, et c'est une limite déclarée. Les faire porter un niveau de
capital serait en revanche indéfendable.
\end{attention}

\paragraph{Apport 1 : structure du risque et mapping DORA.} Sur les 1\,041
incidents, 840 disposent d'un vecteur d'accès initial identifié (81~\%),
regroupés en familles et reliés directement aux articles DORA :

\begin{center}
\renewcommand{\arraystretch}{1.3}
\begin{tabular}{>{\raggedright}p{5cm} c >{\raggedright\arraybackslash}p{6cm}}
\toprule
\rowcolor{lightblue}
\textbf{Vecteur d'attaque} & \textbf{Part} & \textbf{Exigence DORA} \\
\midrule
Phishing / Ingénierie sociale & 38{,}8~\% & Art.~13, sensibilisation \& formation \\
Exploitation vuln.\ exposée & 33{,}8~\% & Art.~26, tests TLPT ; Art.~8--10 \\
Chaîne d'approvisionnement / tiers TIC & 15{,}8~\% & Art.~28--44, gestion du risque tiers \\
Identifiants compromis & 6{,}3~\% & Art.~9, contrôle d'accès \& MFA \\
\bottomrule
\end{tabular}
\end{center}

Deux chiffres sont structurants pour la construction des scénarios du
chapitre~\ref{chap:multietats} : 15{,}8~\% des attaques empruntent le vecteur tiers/supply
chain, coeur des articles~28--44 ; et 49{,}6~\% relèvent de vecteurs
qu'un test TLPT (art.~26) est censé détecter en amont (exploitation et
tiers cumulés).

\paragraph{Limite 1 : biais de visibilité.} Hackmageddon recense des
incidents rendus publics, majoritairement médiatiques ou déclarés, ce qui
sur-représente structurellement les événements les plus visibles au détriment
des incidents ordinaires, moins commentés. La base ne peut donc servir à
estimer un niveau absolu de fréquence, seulement une répartition relative
entre vecteurs d'attaque.

\paragraph{Limite 2 : taxonomie instable.} Une comparaison de structure entre
2023 (3\,019 incidents, Q3 manquant) et 2026 a mis en évidence un artefact de
taxonomie. La catégorie \og Ransomware\fg{} n'existait pas en 2023, les incidents
correspondants étant alors classés sous \og Malware\fg{}. Une lecture naïve
suggérait une hausse du ransomware de +11 points ; après reclassement par
grille de mots-clés commune, le ransomware \emph{recule} en réalité de
23 points.

\begin{center}
\renewcommand{\arraystretch}{1.3}
\begin{tabular}{>{\raggedright}p{4.5cm} c c c}
\toprule
\rowcolor{lightblue}
\textbf{Dimension} & \textbf{2023} & \textbf{2026} & \textbf{$\Delta$} \\
\midrule
Cybercriminalité (motivation) & 81{,}1~\% & 73{,}9~\% & $-7{,}2$ \\
Cyberespionnage (motivation) & 9{,}6~\% & 20{,}6~\% & $+11{,}0$ \\
Ransomware (reclassé) & 35{,}8~\% & 12{,}7~\% & $-23{,}1$ \\
\bottomrule
\end{tabular}
\end{center}

Seules les motivations (taxonomie stable de haut niveau) sont donc
comparables directement dans le temps ; les types d'attaque ne le
sont qu'après un travail de reclassement documenté. Le signal robuste qui
demeure est la bascule du cybercrime de masse vers l'espionnage ciblé, une
information pertinente pour la discussion du risque systémique au chapitre~\ref{chap:resultats},
mais qui illustre aussi la fragilité méthodologique de toute lecture brute
d'une base de ce type.

% SOURCES-SCRIPTS: 63

%--------------------------------------------------------------

%--------------------------------------------------------------

\section{OpRisk Global : la sévérité en euros}
\label{sec:oprisk}

% SOURCES-SCRIPTS: 07 47 63

\paragraph{Nature et périmètre.} OpRisk Global est la base de référence
historique du risque opérationnel au sens de Bâle~2/3, comptant
41\,975 incidents après nettoyage d'une aberration de saisie. Le
risque cyber n'y constitue pas une catégorie native : il est reconstruit à
partir des catégories \og Systems Security\fg{} et \og Business Disruption and
System Failures\fg{}. Le périmètre cyber $\times$ finance retenu compte
582 incidents sur la période 2000--2026. La même base porte l'étude des coûts cyber de
\citet{elingwirfs2019}, qui en extraient de même les pertes d'origine informatique.

\paragraph{Lien avec l'approche AMA.} Cette base est historiquement celle
qui alimentait les modèles internes de type \emph{Advanced Measurement
Approach} évoqués au chapitre~\ref{chap:reglementaire}, avant leur retrait au profit de la SMA sous
Bâle~III/IV. Son usage ici renoue donc, à des fins de modélisation interne et
non de calcul réglementaire de Pilier~1, avec la logique originelle de la LDA
bancaire : calibrer une sévérité extrême à partir de pertes réellement
observées plutôt qu'à partir d'un indicateur de volume forfaitaire.

\paragraph{Statistiques de sévérité.} Elles portent sur la population de la fréquence et de la
descente d'échelle : secteur financier, catégories \textsc{tic}, fenêtre $2005$--$2022$, soit
$1\,476$ observations. La loi de Pareto généralisée, elle, est ajustée sur le périmètre cyber
$\times$ finance du paragraphe précédent, $582$ pertes sur $2000$--$2026$. Les deux populations
ne coïncident pas, et la table ci-dessous décrit la première.

\begin{center}
\renewcommand{\arraystretch}{1.3}
\begin{tabular}{l r l r}
\toprule
\rowcolor{lightblue}
\textbf{Indicateur} & \textbf{Valeur} & \textbf{Indicateur} & \textbf{Valeur} \\
\midrule
Médiane & 1{,}20 M\$ & P99 & 419{,}3 M\$ \\
Moyenne & 33{,}8 M\$ & Maximum & 10\,970{,}8 M\$ \\
P90 & 25{,}3 M\$ & Total cumulé & 49{,}8 Md\$ \\
\bottomrule
\end{tabular}
\end{center}

La concentration extrême de la sévérité constitue une signature de queue lourde qui justifie
directement le recours à la théorie des valeurs extrêmes développée au
chapitre~\ref{chap:cascade} : le décile supérieur porte $92{,}3\,\%$ de la perte totale, et le
seul centile supérieur en porte $65{,}1\,\%$.

% SOURCES-SCRIPTS: 67

\paragraph{Choix du seuil POT.} Le choix du seuil $u$ au-dessus duquel la loi
de Pareto généralisée est calibrée résulte d'un compromis classique
biais-variance, arbitré graphiquement. Le \emph{mean excess plot}
(figure~\ref{don:fig:mep}) présente une pente positive et globalement linéaire
au-delà de $u = 20$~M\euro, signature attendue d'un régime GPD à queue lourde
($\xi > 0$) ; c'est ce seuil, correspondant au percentile~85 de la
distribution des montants, qui est retenu pour la suite.

\begin{figure}[htbp]
\centering
\includegraphics[width=0.85\textwidth]{figures/mean_excess_oprisk.png}
\caption{Mean excess plot sur le périmètre cyber$\times$finance d'OpRisk
Global. La linéarité de $e(u)$ au-delà du seuil retenu $u = 20$~M\euro
(pointillés) confirme le régime de queue lourde justifiant l'ajustement GPD.}
\label{don:fig:mep}
\end{figure}

\paragraph{Calibration GPD et validation croisée.} Au seuil du percentile~85
($u = 20{,}0$~M\euro, 91 excès), la calibration GPD donne :
\[
\hat{\xi} = 0{,}595 \quad (\text{IC}_{90\%}=[0{,}30\,;\,0{,}83]), \qquad
\hat{\sigma} = 58{,}0~\text{M\euro}, \qquad
\qsev_{99{,}5\%}(X) = 663~\text{M\euro}.
\]
Le dernier terme est le quantile de la sévérité d'un sinistre, pas un capital : le
besoin de capital est le quantile de la charge \emph{annuelle agrégée}, obtenu plus loin par
composition avec la fréquence. L'intervalle de confiance bootstrap qui entoure ce quantile est
$[411\,;\,1037]$~M\euro, soit un facteur~2{,}5 entre les bornes, une ampleur qui, à elle seule,
illustre que le SCR ne peut être présenté comme un point.

La comparaison avec la PRC ($\hat{\xi}_{\text{PRC}} = 1{,}03$ contre
$\hat{\xi}_{\text{OpRisk}} = 0{,}595$) ne se lit pas comme une différence de
risque, car les deux indices ne portent pas sur le même objet. La sévérité PRC
n'est pas observée : elle est une puissance du nombre d'enregistrements, de
pente $0{,}76$ en logarithme, et une puissance multiplie l'indice de queue par
son exposant, si bien que $\hat{\xi}_{\text{PRC}}$ décrit la queue des
enregistrements transportée en euros. Le biais de taille d'OpRisk, mesuré faible
au \S\ref{sec:biais-chiffres}, n'en rend pas compte. Les deux sources confirment néanmoins le même régime
qualitatif de queue lourde ($\xi > 0$). Il n'existe donc pas de hiérarchie
\og source primaire\fg{} en sévérité : OpRisk fournit des montants réels, la PRC
un spectre de tailles plus large via la conversion Jacobs : deux
estimations indépendantes et complémentaires, dont l'écart constitue
lui-même une information sur l'incertitude du modèle.

La décomposition par sous-catégorie révèle en outre deux profils de risque
distincts : les pannes (\og Business Disruption\fg{}) sont moins fréquentes
(19~\% des incidents) mais nettement plus coûteuses, avec une médiane cinq
fois supérieure à celle des intrusions (\og Systems Security\fg{}), suggérant que
les multiplicateurs DORA pourraient à terme être différenciés non seulement
en fréquence mais aussi en sévérité, sous réserve de la taille réduite de
l'échantillon \og pannes\fg{} (28 excès), qui rend cette calibration indicative.

%--------------------------------------------------------------

%--------------------------------------------------------------

\section{Portée et limites du dispositif de données}
\label{sec:portee-donnees}

La discipline méthodologique retenue dans ce mémoire consiste à ne jamais
prétendre inférer davantage que ce que les données permettent réellement, ce
qui impose d'expliciter aussi bien les capacités que les limites du dispositif.

\paragraph{Ce que l'on sait faire.} On sait lire une fréquence d'incidents
matériels sur des pertes observées (OpRisk Global), et la moduler entre états de
conformité selon une structure d'attaque citée (Hackmageddon).
On sait calibrer une sévérité de queue lourde par la théorie des valeurs
extrêmes à partir de pertes réellement observées (OpRisk Global), et
confronter cette calibration à une source indépendante pour objectiver
l'incertitude. On sait enfin traiter l'écart entre sources non comme un bruit
à éliminer, mais comme une dimension d'incertitude à part entière, qui
s'ajoute à l'incertitude de paramètre estimée par bootstrap.

\paragraph{Ce que l'on ne sait pas faire.} On ne sait pas mesurer directement
le niveau de conformité DORA d'une entité donnée ; c'est précisément la
raison d'être de l'état de conformité par pilier introduit au
chapitre~\ref{chap:multietats}. On ne sait pas
disposer d'une base d'incidents harmonisée à l'échelle européenne, exhaustive
et non biaisée par la visibilité médiatique. On ne sait pas non plus estimer
avec précision la corrélation entre les différentes briques de perte,
faute d'historique suffisant sur les événements joints ; cette incertitude
est traitée au chapitre~\ref{chap:cascade} par une structure de dépendance explicite plutôt que
par une hypothèse d'indépendance non fondée.

Deux règles du moteur de simulation, le plafond de sévérité propre à la source PRC et la
sévérité nulle sous le seuil, ainsi que la suffisance des trois sources pour un exercice
d'ORSA, sont précisées au \S\ref{sec:ann-sources}.

% SOURCES-SCRIPTS: 05 63

\section{Pourquoi ne pas multiplier les sources}

Des bases
supplémentaires existent (bases commerciales de sinistres cyber assurantiels,
enquêtes sectorielles européennes) et pourraient, en principe, enrichir le
dispositif. Mais chacune porterait ses propres biais de sélection et de
représentativité, sans lever la limite structurelle identifiée plus haut :
l'absence, à ce jour, d'un registre harmonisé et exhaustif à
l'échelle du marché. Multiplier les sources déplacerait la difficulté sans la
résoudre, et complexifierait un dispositif dont les rôles sont aujourd'hui distincts (trois
sources, trois rôles, aucune hiérarchie de \og source primaire\fg{}). L'analyse de sensibilité menée
au chapitre~\ref{chap:resultats} confirme d'ailleurs, \emph{a posteriori}, la pertinence de ce choix :
c'est l'incertitude sur l'indice de queue $\xi$ de la sévérité, estimé sur
les sources déjà mobilisées, qui domine très largement toute autre source
de variation du modèle, multipliers de fréquence et paramètres de dépendance
compris. Approfondir l'exploitation des trois sources retenues (bootstrap,
validation croisée, diagnostics de stabilité) réduit donc l'incertitude du
capital de façon bien plus significative que l'ajout d'une quatrième base de
sinistralité.

%--------------------------------------------------------------

%--------------------------------------------------------------

La structure tirée de Hackmageddon est confrontée au rapport \emph{ENISA Threat Landscape 2025}
\citep{enisa2025}, de méthodologie de collecte distincte, qui confirme la même hiérarchie des
vecteurs d'attaque par leur taux de conversion en intrusion (\S\ref{sec:ann-sources}).

%==============================================================

\begin{cle}
\textbf{Ce que ces données permettent, et ce qu'elles interdisent.} Elles fixent la queue de
sévérité ($\hat\xi = 0{,}595$ sur OpRisk), la fréquence d'entrée et l'accumulation par cause
commune. Elles ne fixent \emph{pas} la direction de la contagion : le
chapitre~\ref{ch:identifiabilite} le montre sur trois sources qui échouent pour des raisons
distinctes, et le
chapitre~\ref{chap:partielle} en tire les conséquences en bornes de capital. Cette frontière
est un résultat produit par le dispositif de données, non une lacune de celui-ci.
\end{cle}

% =====================================================================

% SOURCES-SCRIPTS: 05 62 63
